\chapter{Central Nervous System}
\index{Central Nervous System}

\section*{Definition and Overview}
The central nervous system (CNS) is the control centre of the body. It consists of the brain and the spinal cord, which receive and process information from the body and generate the signals that govern movement, thought, sensation, emotion, and the automatic regulation of vital functions such as breathing and heart rate.

The CNS is protected by the skull and the spine and is wrapped in three layers of membrane (the meninges), bathed in cerebrospinal fluid, and separated from the bloodstream by the blood--brain barrier. The other part of the nervous system, the peripheral nervous system, carries messages to and from the CNS. Damage to the CNS affects every function it controls, which is why neurological disorders are among the most disabling of all diseases.

\section*{Epidemiology}
Disorders of the central nervous system are among the most important causes of disability and death worldwide. Stroke is a leading cause of adult disability, and dementia, particularly Alzheimer's disease, affects a growing number of older people as populations age.

Epilepsy affects roughly 1 in 100 people, while conditions such as multiple sclerosis, Parkinson's disease, and motor neurone disease are less common but have a major impact on those affected. Head injuries and infections of the nervous system, such as meningitis, can affect people of any age and are medical emergencies.

\section*{Aetiology and Pathophysiology}
CNS disorders can arise from many different processes:

\begin{itemize}
    \item \textbf{Vascular:} stroke from a blocked or bleeding vessel deprives brain tissue of blood and oxygen
    \item \textbf{Degenerative:} progressive loss of nerve cells, as in Alzheimer's and Parkinson's disease
    \item \textbf{Inflammatory and autoimmune:} the immune system attacks the nervous system, as in multiple sclerosis
    \item \textbf{Infectious:} meningitis, encephalitis, and brain abscesses caused by viruses, bacteria, or fungi
    \item \textbf{Traumatic:} concussion, brain injury, and spinal cord injury from impact
    \item \textbf{Neoplastic:} primary brain tumours or cancers spreading to the brain from elsewhere
    \item \textbf{Epileptic:} abnormal electrical activity producing seizures
    \item \textbf{Developmental and genetic:} conditions present from birth or inherited, such as cerebral palsy or Huntington's disease
\end{itemize}

Nerve cells (neurons) transmit signals rapidly along pathways; damage to these cells or their connections interrupts the functions they control, and, unlike many other tissues, the CNS has very limited ability to regenerate.

\section*{Clinical Features}
Symptoms depend on which part of the nervous system is affected:

\begin{itemize}
    \item Weakness, paralysis, or loss of coordination on one side of the body
    \item Numbness, tingling, or altered sensation
    \item Difficulty speaking, understanding, or swallowing
    \item Vision problems, including loss of vision or double vision
    \item Headache, often severe or persistent
    \item Seizures, blackouts, or loss of consciousness
    \item Confusion, memory loss, or changes in behaviour and personality
    \item Loss of bladder or bowel control
    \item Dizziness, unsteadiness, or problems with balance
\end{itemize}

\section*{Differential Diagnosis}
\begin{itemize}
    \item Peripheral nervous system disorders, such as neuropathy or nerve entrapment
    \item Muscle disorders, such as muscular dystrophy
    \item Psychiatric conditions that can mimic neurological symptoms
    \item Migraine and other primary headache syndromes
    \item Metabolic causes of confusion, such as low blood sugar or electrolyte disturbance
    \item Infections and autoimmune conditions affecting other organs that produce generalised symptoms
\end{itemize}

\section*{When to Seek Medical Care}
See a doctor for persistent headaches, new weakness or sensory change, memory problems, or recurrent seizures or blackouts.

\textbf{Contact your local emergency services for:}
\begin{itemize}
    \item Sudden weakness or drooping of the face, arm, or leg, or difficulty speaking (possible stroke; think FAST: Face, Arms, Speech, Time)
    \item Sudden severe headache, often described as the worst ever
    \item A seizure lasting more than five minutes, or repeated seizures
    \item Loss of consciousness or a sudden collapse
    \item Neck stiffness with fever, rash, or severe headache (possible meningitis)
    \item Head injury with confusion, vomiting, or drowsiness
\end{itemize}

\section*{Diagnosis and Investigations}
\begin{itemize}
    \item \textbf{History and examination:} careful neurological examination of strength, sensation, coordination, reflexes, and eye movements
    \item \textbf{Imaging:} CT scan for rapid assessment of bleeding or stroke; MRI for detailed images of the brain and spinal cord
    \item \textbf{Electroencephalogram (EEG):} records electrical activity and is used in epilepsy
    \item \textbf{Lumbar puncture:} samples cerebrospinal fluid to diagnose infection, inflammation, or bleeding
    \item \textbf{Nerve conduction studies:} distinguish CNS from peripheral nerve problems
    \item \textbf{Blood tests:} infection markers, electrolytes, vitamin levels, and autoimmune antibodies
    \item \textbf{Cognitive testing:} structured assessment of memory and thinking
\end{itemize}

\section*{Management}
\begin{itemize}
    \item \textbf{Acute emergencies:} clot-busting or clot-removal treatment for stroke, and urgent treatment for meningitis or severe head injury
    \item \textbf{Medication:} anti-seizure drugs, medicines for Parkinson's disease, immunotherapies for multiple sclerosis, and treatment of the underlying cause
    \item \textbf{Rehabilitation:} physiotherapy, occupational therapy, and speech therapy to restore function after stroke or injury
    \item \textbf{Supportive care:} management of bladder, swallowing, and mobility problems
    \item \textbf{Surgery:} removal of tumours, relief of pressure on the brain or spinal cord, or treatment of bleeding
    \item \textbf{Psychological support:} counselling and cognitive rehabilitation for memory, mood, and behaviour
    \item \textbf{Long-term monitoring:} regular specialist review for progressive conditions
\end{itemize}

\section*{Complications}
\begin{itemize}
    \item Permanent disability, including paralysis, speech, and swallowing problems
    \item Cognitive decline, memory loss, and changes in personality
    \item Epilepsy following stroke, head injury, or infection
    \item Pressure sores, chest infections, and blood clots from immobility
    \item Loss of independence and need for long-term care
    \item Falls and fractures related to unsteadiness or weakness
    \item Depression and anxiety, which are common after neurological illness
\end{itemize}

\section*{Prognosis and Outcomes}
Outcomes vary widely depending on the condition and how quickly treatment begins. Stroke and head injury outcomes are strongly influenced by early treatment and rehabilitation, and many people recover substantially in the months after the event. Some conditions, such as Parkinson's disease and dementia, are progressive and managed over the long term, while others, including many epilepsies, can be well controlled. Rehabilitation plays a central role in improving function and quality of life, and many people live full lives with appropriate support.

\section*{Prevention and Self-Care}
\begin{itemize}
    \item Control blood pressure, cholesterol, and diabetes to reduce stroke risk
    \item Do not smoke; limit alcohol
    \item Stay physically and mentally active
    \item Wear helmets and seatbelts to prevent head and spine injury
    \item Maintain up-to-date vaccinations against meningitis-causing bacteria
    \item Take prescribed medicines for seizures or other conditions reliably
    \item Protect sleep and manage stress, which support overall brain health
    \item Follow rehabilitation programs and attend specialist reviews
\end{itemize}

\section*{Sources and Further Reading}
The following organisations provide reliable information on the central nervous system and neurological disorders.

\begin{itemize}
    \item National Health Service. \textit{NHS conditions guide on brain, nerve, and spine conditions.} https://www.nhs.uk/conditions/
    \item Mayo Clinic. \textit{Neurological disease information and patient education.} https://www.mayoclinic.org/
    \item World Health Organization. \textit{Neurological disorders: global burden and policy.} https://www.who.int/
    \item MedlinePlus. \textit{Consumer health information on brain and nervous system disorders.} https://medlineplus.gov/
    \item MSD Manual. \textit{Neurologic disorders -- professional overview.} https://www.msdmanuals.com/
    \item Centers for Disease Control and Prevention. \textit{Brain health and stroke information.} https://www.cdc.gov/
\end{itemize}
